\chapter{Introducción}
\label{sec:introduccion}

\section{Antecedentes}
\label{sec:intro-antecedentes}

El Centro Óptico Santa María gestionaba su operación --- agendamiento de citas, historia clínica de pacientes, control de inventario y seguimiento de ventas --- de forma manual y fragmentada, sin una plataforma que centralizara esa información. Sobre esa necesidad, en el marco de la Licenciatura en Ciencias Informáticas de la Facultad Politécnica de la Universidad Nacional de Asunción, el anteproyecto relevó el problema y propuso el desarrollo de SIGA-Óptica: un sistema integral de gestión y agendamiento que unifique las operaciones administrativas, clínicas y comerciales de la óptica y su consultorio oftalmológico asociado.

El anteproyecto definió el problema, los objetivos, el modelo de casos de uso de alto nivel y un diseño preliminar de entidades. Este documento no repite ese análisis; lo da por sentado y se concentra en reportar lo que efectivamente se construyó.

\section{Objetivo del documento}
\label{sec:intro-objetivo-documento}

Este documento constituye el informe del trabajo final de grado y reporta el resultado del desarrollo de SIGA-Óptica: la arquitectura de la solución finalmente adoptada, la construcción de cada módulo de negocio, el modelo de datos, el proceso de implementación y despliegue, la validación realizada sobre el sistema y los cambios introducidos respecto al diseño planteado en el anteproyecto, con su justificación.

\section{Alcance}
\label{sec:intro-alcance}

Quedan dentro del sistema implementado y reportado en este documento los dieciséis módulos de negocio detallados a partir del Capítulo~\ref{mod:01-identidad-y-acceso}: identidad y acceso, pacientes y profesionales, personal, agenda de turnos, historia clínica, clientes, inventario y catálogo, ventas, laboratorio óptico, compras, caja y egresos, roles y permisos, notificaciones, reportes, soporte multi-sucursal y auditoría.

Quedan fuera de este entregable --- y se documentan como tales en el Capítulo~\ref{sec:cambios-proyecto1} y en las Conclusiones (Capítulo~\ref{sec:conclusiones}) --- las funcionalidades que el anteproyecto contemplaba y no se llegaron a implementar en esta etapa, así como las limitaciones de diseño conocidas del sistema actual.
